

% \subsection{Financial metrics}

% Average slippage, the average absolute difference between the executed price and the intended price. 

% \begin{equation}
%     S_\text{avg} = \frac{1}{N}\sum_{i=1}^{N}\mid P_{\text{exec},i} - P_{\text{int},i}  \mid 
% \end{equation}

% Fraction of the requested trading volume that was actually executed.

% \begin{equation}
%     \text{Fill Rate} = \frac{\sum_{i=1}^{N} V_{\text{exec},i}}{\sum_{i=1}^{N}V_{\text{req},i}}
% \end{equation}

% The average time a position remains open.

% \begin{equation}
%     \overline{T} = \frac{1}{N}\sum_{i=1}^{N}\left( t_{\text{close},i} - t_\text{open},i \right)
% \end{equation}

% The maximum drawdown, largest peak-to-through decline in portfolio value over the period

% \begin{equation}
%     \text{MDD} = \max_{\tau \in (0, T)} \left( \max_{t \in (0, \tau)} \frac{V(t) - V(\tau)}{V(t)} \right)
% \end{equation}

% The Sharpe ratio, excess return per unit of total volatility.

% \begin{equation}
%     \text{Sharpe} = \frac{\mathbb{E}[R_p - R_f]}{\sigma_p}
% \end{equation}

% The Sortino ratio, the average excess return per unit of downside volatility.

% \begin{equation}
%     \text{Sortino} = \frac{E[R_p - R_f]}{\sigma_d}
% \end{equation}

% The win rate, the fraction of trades that are profitable. 

% \begin{equation}
%     W = \frac{N_{\text{win}}}{N_{\text{total}}}
% \end{equation}

% Profit factor, ratio of total profit to total losses.

% \begin{equation}
%     \text{PF} = \frac{\sum \text{Gross Profits}}{\sum |\text{Gross Losses}|}
% \end{equation}

% Average net profit er trade after transaction costs.

% \begin{equation}
%     \text{PnL}_{\text{avg}} = \frac{1}{N} \sum_{i=1}^{N} \left( (P_{\text{exit}, i} - P_{\text{entry}, i})Q_i - c_i\right)
% \end{equation}

% Cumulative net profit and loss, the total economic value of the strategy, rather than just per-trade averages.

% \begin{equation}
%     \text{PnL}_\text{tot} = \sum_{i=1}^{N}\left((P_{\text{exit},i}-P_{\text{entry},i})Q_i-c_i\right)
% \end{equation}

% Return volatility, while Sharpe and Sortino report volatility implicitly, this reports it directly. 

% \begin{equation}
%     \sigma_R=\sqrt{\frac{1}{T-1}\sum_{t=i}^{T}\left(R_t-\bar{R}\right)^2}
% \end{equation}
% \begin{equation}
%     \bar{R}=\frac{1}{T}\sum_{t=1}^{T}R_t
% \end{equation}

% Calmar ratio, could be more informative than Sharpe for trading when drawdown matters. 

% \begin{equation}
%     \text{Calmar} = \frac{\mathbb{E}[R_p-R_f]}{\text{MDD}}
% \end{equation}
% \begin{equation}
%     \text{Calmar} = \frac{R_\text{ann}}{\text{MDD}}
% \end{equation}

% Average trade per trade as a fraction of capital committed.

% \begin{equation}
%     \bar{r}_\text{trade} = \frac{1}{N}\sum_{i=1}^{N}\frac{\left( P_{\text{exit},i} - P_{\text{entry},i}\right)Q_i-c_i}{P_{\text{entry},i}Q_i}
% \end{equation}

% Turnover, measure of how intensively the strategy trades relative to available capital, where $A_0$ is initial capital.

% \begin{equation}
%     \text{Turnover} = \frac{\sum_{i=1}^{N}\left|Q_i\right|P_{\text{entry},i}}{A_0}
% \end{equation}

% Exposure time, the fraction of total sample time during which the strategy has an open position.

% \begin{equation}
%     \text{Exposure} = \frac{\sum_{i=1}^{N}\left( t_{\text{close},i} - t_{\text{open},i} \right)}{T_\text{sample}}
% \end{equation}

% The average loss and average gain, separates how large winning trades are from how large losing trades are. 

% \begin{equation}
%     \overline{\text{Gain}} = \frac{1}{N_\text{win}}\sum_{i\in\mathcal{W}}\text{PnL}_i
% \end{equation}
% \begin{equation}
%     \overline{\text{Loss}} = \frac{1}{N_\text{loss}}\sum_{i\in\mathcal{L}}\left|\text{PnL}_i\right|
% \end{equation}

% Expectancy, the expected profit per trade, combining both win rates and win/loss magnitude.

% \begin{equation}
%     \text{Expectancy} = W\cdot\overline{\text{Gain}}-(1-W)\cdot\overline{\text{Loss}}
% \end{equation}

